\documentclass[11pt,a4paper]{article}

% ── Packages ──────────────────────────────────────────────────────────
\usepackage[utf8]{inputenc}
\usepackage[T1]{fontenc}
\usepackage{lmodern}
\usepackage{microtype}
\usepackage{amsmath,amssymb}
\usepackage{booktabs}
\usepackage{enumitem}
\usepackage{hyperref}
\usepackage{listings}
\usepackage{xcolor}
\usepackage{graphicx}
\usepackage{float}
\usepackage[margin=2cm]{geometry}
\usepackage{caption}
\usepackage{subcaption}
\usepackage{algorithm}
\usepackage{algpseudocode}
\usepackage{fancyhdr}
\usepackage{titlesec}
\usepackage{parskip}

% ── Hyperref Config ──────────────────────────────────────────────────
\hypersetup{
  colorlinks=true,
  linkcolor=blue!70!black,
  citecolor=green!50!black,
  urlcolor=blue!60!black,
  pdftitle={RAGLeakLab: Deterministic Security Testing for RAG Systems},
  pdfauthor={RAGLeakLab Contributors},
}

% ── Listings Config ──────────────────────────────────────────────────
\definecolor{codebg}{RGB}{248,248,248}
\definecolor{codeframe}{RGB}{200,200,200}
\definecolor{codekw}{RGB}{0,112,32}
\definecolor{codecomment}{RGB}{96,96,96}
\definecolor{codestr}{RGB}{186,33,33}

\lstset{
  backgroundcolor=\color{codebg},
  frame=single,
  rulecolor=\color{codeframe},
  basicstyle=\ttfamily\scriptsize,
  keywordstyle=\color{codekw}\bfseries,
  commentstyle=\color{codecomment}\itshape,
  stringstyle=\color{codestr},
  breaklines=true,
  breakatwhitespace=true,
  numbers=left,
  numberstyle=\tiny\color{gray},
  tabsize=2,
  showstringspaces=false,
  language=Python,
}

% ── Title ────────────────────────────────────────────────────────────
\title{%
  \textbf{RAGLeakLab: Deterministic Security Testing\\
  for Retrieval-Augmented Generation Systems}
}
\author{%
  RAGLeakLab Contributors\\
  \texttt{https://github.com/mishabar410/RAGLeakLab}
}
\date{February 2026 \quad --- \quad v1.0.0}

\begin{document}
\maketitle

% ══════════════════════════════════════════════════════════════════════
% ABSTRACT
% ══════════════════════════════════════════════════════════════════════
\begin{abstract}
Retrieval-Augmented Generation (RAG) systems combine document retrieval with
language model generation, introducing a dual attack surface: the corpus
may leak through retrieval and the generator may reproduce sensitive content.
Existing security evaluations rely on stochastic, model-dependent probes that
produce non-reproducible results, making them unsuitable for continuous
integration and compliance workflows.

We present \textsc{RAGLeakLab}, an open-source framework that addresses
these limitations through a principled, deterministic design. The framework
implements six threat classes---canary extraction, verbatim reproduction,
membership inference, semantic leakage, cross-document correlation, and
corpus poisoning---evaluated via metrics that require no neural inference at
test time. A TF-IDF reference pipeline, fixed paraphrase templates, and
deterministic delta debugging ensure identical results across runs. CI
integration is provided through JUnit XML and SARIF export, HTTP
record/replay cassettes for network-free execution, and configurable
threshold calibration with target false-positive rate constraints.

RAGLeakLab is implemented in approximately 12,000 lines of Python, carries
no GPU dependencies, and produces bit-identical reports across environments.
The framework is designed for integration into enterprise CI/CD pipelines
where reproducibility, auditability, and offline operation are requirements.
\end{abstract}

% ══════════════════════════════════════════════════════════════════════
\section{Introduction}\label{sec:intro}
% ══════════════════════════════════════════════════════════════════════

Retrieval-Augmented Generation (RAG) has become the dominant architecture for
grounding large language models in private document collections. By retrieving
relevant passages and injecting them into the generation context, RAG systems
reduce hallucination and enable knowledge-intensive applications without fine-tuning.
However, this architecture introduces a fundamental tension: the same retrieval
mechanism that grounds generation also creates pathways for information leakage.

An attacker who can observe the generated output---and in many deployments, any
authenticated user qualifies---may extract sensitive content from the private
corpus through carefully crafted queries. The threat surface spans multiple
dimensions: verbatim reproduction of source text, inference of corpus membership,
extraction of planted canary tokens, semantic exposure of sensitive facts,
cross-document correlation of information that should not be composed, and
integrity compromise through corpus poisoning.

Evaluating these threats in practice presents a reproducibility challenge.
Existing approaches often rely on neural classifiers, stochastic probing, or
model-specific adversarial techniques that produce different results across
runs, hardware, and software versions. Such non-determinism is fundamentally
incompatible with CI/CD pipelines, regression testing, and compliance
auditing, where a security gate must return the same verdict given the same
inputs.

\textsc{RAGLeakLab} is designed from first principles to be deterministic.
Every component---from query transformation to metric calculation to report
generation---produces identical output across runs when given identical input.
The framework achieves this through several design choices:

\begin{itemize}[nosep]
  \item A TF-IDF reference retriever with deterministic tie-breaking replaces
        embedding-model-dependent retrieval.
  \item Fixed template-based paraphrase generation replaces neural paraphrasing.
  \item Longest Common Substring and n-gram overlap replace neural similarity
        for verbatim detection.
  \item Pattern-matching claim detectors replace LLM-based fact verification.
  \item Delta debugging for query minimization uses a deterministic algorithm.
  \item SHA-256 provenance hashes record corpus, attack, and configuration
        state for auditability.
\end{itemize}

The remainder of this paper is structured as follows. Section~\ref{sec:problem}
defines the problem space. Section~\ref{sec:threat} presents the formal threat
model. Section~\ref{sec:design} describes design principles.
Section~\ref{sec:arch} details the system architecture.
Section~\ref{sec:attacks} covers the attack framework.
Section~\ref{sec:metrics} formalizes the leakage and integrity metrics.
Section~\ref{sec:poisoning} describes dataset poisoning evaluation.
Section~\ref{sec:calibration} covers threshold calibration.
Section~\ref{sec:attribution} presents the attribution engine.
Section~\ref{sec:ci} describes CI integration.
Section~\ref{sec:benchmarks} provides benchmark results.
Section~\ref{sec:security} discusses security hardening.
Section~\ref{sec:limitations} acknowledges limitations.
Section~\ref{sec:related} surveys related work, and
Section~\ref{sec:conclusion} concludes.

% ══════════════════════════════════════════════════════════════════════
\section{Problem Statement}\label{sec:problem}
% ══════════════════════════════════════════════════════════════════════

A production RAG system serves as an interface between untrusted user queries
and a private document corpus. The security requirement is that the system
should answer questions using corpus knowledge without exposing sensitive
content beyond what is strictly necessary for the answer. Formally, let
$\mathcal{C} = \{d_1, \ldots, d_n\}$ be a private corpus, $q$ a user query,
$R(q, \mathcal{C}) \subseteq \mathcal{C}$ the retrieved subset, and
$G(q, R)$ the generated answer. The security testing problem is to
determine whether $G(q, R)$ reveals information from $\mathcal{C}$ that
constitutes a confidentiality violation.

Three properties distinguish this problem from general LLM safety testing:

\textbf{Dual attack surface.} Both the retriever $R$ and the generator $G$
contribute to leakage. The retriever may return documents containing secrets;
the generator may verbatim reproduce, paraphrase, or compose information
from the retrieved context.

\textbf{Corpus-dependent threats.} The severity of leakage depends on the
specific corpus content. A canary token \texttt{CANARY\_<hex>} planted in the
corpus constitutes a definitive leak when reproduced; verbatim overlap with
public boilerplate does not. Testing must be parameterized by corpus content
and associated claim annotations.

\textbf{Reproducibility requirements.} Enterprise deployments require
security gates that produce identical verdicts given identical inputs, support
regression detection across releases, and generate audit trails for compliance.
Non-deterministic probes violate all three requirements.

RAGLeakLab addresses this problem by providing a testing framework where
every component is deterministic, every threat class has a well-defined metric
with a configurable threshold, and every run produces a versioned,
hash-anchored report suitable for CI gates.

% ══════════════════════════════════════════════════════════════════════
\section{Threat Model}\label{sec:threat}
% ══════════════════════════════════════════════════════════════════════

RAGLeakLab implements six threat classes, each corresponding to a distinct
information leakage vector. The framework explicitly separates \emph{leakage
threats} (Sections~\ref{subsec:canary}--\ref{subsec:crossdoc}), which
concern information escaping from the corpus, from \emph{integrity threats}
(Section~\ref{subsec:integrity}), which concern adversarial manipulation
of the corpus itself.

\subsection{Canary Extraction}\label{subsec:canary}

Canary tokens are deterministic, high-entropy strings of the form
\texttt{CANARY\_<16 hex chars>}, generated via SHA-256 hashing of a seed:
\texttt{SHA256("ragleaklab\_canary\_\{seed\}")[:16]}. Canaries are injected
into corpus documents using configurable strategies (end-of-document or
mid-document insertion). Their presence in generated output constitutes an
unambiguous confidentiality violation---the token has zero probability of
appearing in any legitimate response.

\textbf{Attacker model.} The attacker submits queries designed to elicit
reproduction of injected secrets. Strategies include direct requests,
indirect framing, role confusion, and context manipulation
(Section~\ref{sec:attacks}).

\textbf{Detection.} A compiled regular expression
\texttt{CANARY\_[0-9a-f]\{16,\}} is applied to the generated output. The
metric reports count and presence as a binary signal.

\subsection{Verbatim Reproduction}\label{subsec:verbatim}

Verbatim leakage measures the degree of exact textual overlap between the
generated output and source documents. Two complementary metrics are computed:

\begin{itemize}[nosep]
  \item \textbf{Longest Common Substring (LCS) ratio}: The length of the
        longest common substring between the answer and any source document,
        normalized by source document length. A minimum length threshold
        (default: 20 characters) suppresses common-phrase false positives.
  \item \textbf{N-gram overlap}: The count of matching n-grams (default:
        $n = 5$) between answer and source, with a configurable minimum
        match count.
\end{itemize}

The combined score is $\max(\text{lcs\_ratio}, \text{ngram\_ratio})$, where
ngram\_ratio is the fraction of source n-grams appearing in the answer.

\subsection{Membership Inference}\label{subsec:membership}

Membership inference determines whether a specific document was present in
the corpus based on observable response characteristics. RAGLeakLab implements
two approaches:

\textbf{Heuristic scoring.} A weighted combination of retrieval score
statistics (mean, max), answer length, and context specificity indicators.
When non-member reference signals are available, the score is the normalized
difference between member and non-member signal strengths.

\textbf{Paraphrase consistency.} The query is expanded into deterministic
paraphrases using fixed templates (e.g., ``What is \{query\}?'',
``Tell me about \{query\}'', ``Details on \{query\}''). Retrieved document
overlap and answer token similarity across paraphrases are measured.
High consistency suggests corpus membership. The consistency score combines
retrieval stability (Jaccard index over retrieved document IDs) and answer
similarity (token-level Jaccard).

\subsection{Semantic Leakage}\label{subsec:semantic}

Semantic leakage evaluates whether sensitive \emph{claims}---discrete facts
annotated in the corpus---are revealed in the generated output. Claims are
defined per-document in JSONL format with typed sensitivity levels (high,
medium, low) and categorical types (e.g., \texttt{EMAIL}, \texttt{PHONE},
\texttt{ACCOUNT\_ID}, \texttt{SECRET\_CODEWORD}, \texttt{INTERNAL\_PROJECT}).

Detection uses type-specific matchers:
\begin{itemize}[nosep]
  \item \textbf{Email}: Regex-based extraction and exact matching.
  \item \textbf{Phone}: E.164 and common format normalization, then
        digit-sequence comparison.
  \item \textbf{Account ID / Secret codeword}: Case-insensitive substring
        matching with text normalization (whitespace collapse, punctuation
        removal).
  \item \textbf{General}: Normalized substring containment after text
        preprocessing.
\end{itemize}

The semantic leakage rate is $|\text{leaked claims}| / |\text{total claims
in retrieved documents}|$.

\subsection{Cross-Document Correlation}\label{subsec:crossdoc}

Cross-document leakage detects composed claims---facts that require
information from \emph{multiple} source documents to construct. These are
defined as \emph{composed claims} in a separate JSONL file, each referencing
two or more source document IDs. Detection verifies that a matched claim
indeed spans multiple distinct document origins, preventing single-document
matches from triggering false positives.

The metric reports the count of leaked composed claims and provides evidence
linking each match to its source documents.

\subsection{Corpus Poisoning (Integrity Threats)}\label{subsec:integrity}

Integrity threats model an adversary who can inject or modify documents in
the corpus. RAGLeakLab implements three sub-categories, each with a
dedicated evidence schema:

\textbf{Retrieval hijacking} (\texttt{RetrievalIntegrityEvidence}).
Poisoned documents are injected to rank higher than legitimate content for
targeted queries. Detection compares expected vs.\ actual document IDs in
retrieval results, computing poison rate at $k$
($|\text{poison docs in top-}k| / k$), poison mean reciprocal rank, and
true document recall at $k$.

\textbf{Claim corruption} (\texttt{ClaimIntegrityEvidence}). Injected
documents contain false claims designed to appear in generated output.
Detection matches expected true claims against actual output and identifies
the presence of poison claims, reporting contradiction counts and
corruption rates.

\textbf{Sentinel takeover} (\texttt{SentinelIntegrityEvidence}). Backdoor
triggers planted in the corpus activate when specific patterns appear in
queries. Detection uses rule-based pattern matching with deterministic
policies (block/strip/allow) and reports whether triggers were activated,
the policy action taken, and output markers found.

Each evidence type carries a severity level (high/medium/low), confidence
score, pack and query identifiers, and optional detail dictionaries.

% ══════════════════════════════════════════════════════════════════════
\section{Design Principles}\label{sec:design}
% ══════════════════════════════════════════════════════════════════════

Five design principles govern RAGLeakLab's architecture:

\subsection{Determinism}

Every operation produces identical output given identical input. The
reference pipeline uses a TF-IDF retriever with deterministic tie-breaking
(secondary sort by chunk ID ensures stable ranking when cosine similarities
are equal). Paraphrase generation uses fixed templates rather than neural
models. Delta debugging follows a deterministic splitting algorithm.
Report entries are sorted by \texttt{test\_id} before serialization. A
dedicated verification module (\texttt{core.determinism}) can run the same
pack $N$ times and compare normalized outputs, stripping volatile fields
(timestamps, timings) before deep comparison.

\subsection{Provenance}

Every report carries SHA-256 hashes of its inputs: corpus directory content
(\texttt{corpus\_hash}), attack definitions (\texttt{attacks\_hash}),
runtime configuration (\texttt{config\_hash}), and target system identifier
(\texttt{target\_hash}). The \texttt{compute\_config\_hash} function
deterministically serializes key-value settings, sorts them, and hashes the
result. This allows auditors to verify that two reports were produced from
identical inputs without accessing the inputs themselves.

\subsection{Separation of Concerns}

The framework cleanly separates five subsystems: (1)~the attack framework
(case loading, strategy application, execution), (2)~target adapters
(in-process pipeline, HTTP endpoints), (3)~metric calculators (canary,
verbatim, membership, semantic, crossdoc, integrity), (4)~the reporting
layer (schema, export, regression diff), and (5)~corpus management
(loading, patching, claim annotation, canary injection). Each subsystem
communicates through typed Pydantic contracts defined in
\texttt{core.contracts}.

\subsection{Extensibility}

A plugin registry (\texttt{core.plugins}) supports three extension points:
metrics, attacks, and targets. Plugins can be registered programmatically
via \texttt{register(kind, name, obj)} or discovered automatically through
Python entry points in the \texttt{ragleaklab.\{kind\}} namespace. Built-in
components self-register during module import, ensuring they are available
without explicit configuration.

\subsection{CI-First Design}

Every feature is designed for unattended execution. HTTP cassettes enable
network-free replays. JUnit XML and SARIF exports integrate with standard
CI tooling. Exit codes are structured (0=success, 1=general error,
2=input error, 3=validation error, 4=config error, 5=network error,
99=internal error). Error messages are automatically sanitized via
\texttt{core.errors} to prevent secret leakage in CI logs.

% ══════════════════════════════════════════════════════════════════════
\section{System Architecture}\label{sec:arch}
% ══════════════════════════════════════════════════════════════════════

\subsection{Reference RAG Pipeline}

RAGLeakLab includes a reference RAG pipeline (\texttt{rag.pipeline.RAGPipeline})
composed of three stages:

\textbf{Retrieval.} The \texttt{TFIDFRetriever} chunks documents using
configurable window and overlap parameters (default: 256 characters, 32
overlap), builds TF-IDF vectors with $\text{TF} = \log(1 + \text{count})$
and $\text{IDF} = \log(N / \text{df})$, and ranks by cosine similarity.
Ties are broken by chunk ID (lexicographic, ascending) for stable ordering.

\textbf{Context assembly.} The \texttt{ContextBuilder} concatenates
retrieved chunks into a single context string passed to the generator.

\textbf{Generation.} The \texttt{MockGenerator} produces deterministic
answers by echoing retrieved context, ensuring the pipeline's output is
fully determined by its inputs without requiring any language model.

This reference pipeline serves as a controlled baseline. Production systems
are tested via the \texttt{Target} adapter protocol.

\subsection{Target Adapters}

The \texttt{Target} protocol defines a single method:
\texttt{ask(query) -> TargetResponse}, where \texttt{TargetResponse}
contains the answer text, optional retrieved document IDs, optional
retrieval scores, and raw response data. Two implementations are provided:

\textbf{InProcessTarget.} Wraps a \texttt{RAGPipeline} instance, enabling
direct testing without network overhead.

\textbf{HttpTarget.} Sends queries to external RAG endpoints via HTTP.
Configuration includes request method, URL, field mapping for request/response
payloads, custom headers, timeout, and rate limiting (configurable
requests-per-second with a token bucket). Security features include SSRF
protection (Section~\ref{sec:security}) and domain allowlisting.

The \texttt{HttpTarget} supports three operational modes:
\begin{itemize}[nosep]
  \item \texttt{live}: Normal HTTP requests.
  \item \texttt{record}: Requests are executed and responses are recorded
        to JSONL cassette files.
  \item \texttt{replay}: Responses are served from cassette files without
        network access.
\end{itemize}

\subsection{Contracts}

All inter-module communication uses typed Pydantic models defined in
\texttt{core.contracts}:

\begin{itemize}[nosep]
  \item \texttt{Document}, \texttt{Chunk}, \texttt{RetrievalHit}: Corpus
        and retrieval data structures.
  \item \texttt{RunArtifact}: Complete result of executing a single attack
        case, including query, answer, retrieved chunks, timings, context
        statistics, provenance hashes, and optional integrity evidence.
  \item \texttt{MetricScore}: Unified metric result with name, value,
        details dictionary, and optional pass/fail status.
  \item \texttt{CaseResult}: Combines a \texttt{RunArtifact} with its
        associated \texttt{MetricScore} list and overall verdict.
  \item \texttt{ReportSummary}: Top-level report with schema version,
        tool version, timestamp, overall pass/fail, aggregate metrics,
        failure details, and metadata.
\end{itemize}

\subsection{Disk Cache}

The \texttt{core.cache.DiskCache} stores pipeline results keyed by a
SHA-256 hash of corpus hash, target hash, query text, and retrieval
parameters. This enables fast repeat runs when inputs are unchanged. Cache
entries are stored as individual JSON files, using the key hash as filename.
Invalid entries are automatically removed on read.

\subsection{Corpus Management}

The corpus subsystem provides:
\begin{itemize}[nosep]
  \item \textbf{Document loading}: Plain text files from a directory.
  \item \textbf{Claim annotation}: JSONL files mapping \texttt{(doc\_id,
        claim\_id)} to sensitive facts with type and sensitivity metadata.
  \item \textbf{Canary injection}: Deterministic token generation from
        integer seeds and configurable injection strategies (end/middle).
  \item \textbf{Corpus patching}: Deterministic application of incremental
        changes (add/replace/remove documents and claims) via YAML/JSON
        patch specifications (\texttt{corpus.patch}), producing
        hash-verified output corpora.
\end{itemize}

% ══════════════════════════════════════════════════════════════════════
\section{Attack Framework}\label{sec:attacks}
% ══════════════════════════════════════════════════════════════════════

\subsection{Attack Case Schema}

Each attack is defined as an \texttt{AttackCase} with the following fields:
\texttt{test\_id} (unique identifier), \texttt{threat} (threat class),
\texttt{strategy} (transformation to apply), \texttt{query} or
\texttt{turns} (single-turn string or multi-turn conversation), optional
\texttt{expected} output, and optional \texttt{tags} for filtering.
Pydantic validation enforces that exactly one of \texttt{query} or
\texttt{turns} is provided. Multi-turn attacks are flattened to a single
effective query via the \texttt{effective\_query} property, which
concatenates turn roles and content.

Attack cases are loaded from YAML files organized into \emph{packs}---named
collections of test cases targeting specific threat classes.

\subsection{Strategy Catalog}

The \texttt{attacks.catalog} module maintains a registry of transformation
strategies, each defined as an \texttt{AttackStrategy} dataclass with
name, description, and a \texttt{transform: Callable[[str], str]} function.
Built-in strategies include:

\begin{table}[H]
\centering
\caption{Built-in attack strategies.}
\label{tab:strategies}
\small
\begin{tabular}{@{}ll@{}}
\toprule
\textbf{Strategy} & \textbf{Transformation} \\
\midrule
\texttt{direct\_ask} & Identity (no transformation) \\
\texttt{indirect\_ask} & Indirect framing \\
\texttt{role\_confusion} & System role manipulation \\
\texttt{context\_manipulation} & Context injection \\
\texttt{semantic\_leakage} & Semantic probing \\
\texttt{crossdoc\_correlation} & Cross-document composition \\
\texttt{authority\_appeal} & Authority framing \\
\texttt{hypothetical\_framing} & Hypothetical scenarios \\
\texttt{metadata\_probe} & Metadata extraction \\
\texttt{enumeration} & Systematic enumeration \\
\bottomrule
\end{tabular}
\end{table}

Custom strategies can be added via the plugin registry.

\subsection{Execution Engine}

The \texttt{attacks.runner} module orchestrates test execution:

\begin{enumerate}[nosep]
  \item \textbf{Case loading}: \texttt{load\_cases(pack\_path)} parses YAML
        files into \texttt{AttackCase} instances.
  \item \textbf{Strategy application}: The named strategy's
        \texttt{transform} function modifies the query.
  \item \textbf{Cache lookup}: If a \texttt{DiskCache} is provided, the
        cache key (corpus hash, target hash, transformed query, retrieval
        parameters) is checked.
  \item \textbf{Execution}: The query is sent through either a
        \texttt{RAGPipeline} or a \texttt{Target} adapter.
  \item \textbf{Artifact creation}: A \texttt{RunArtifact} is constructed
        with the answer, retrieved chunks, context, timing breakdown
        (retrieval, generation, total milliseconds), context statistics
        (character count, chunk count, truncation flag), and provenance
        hashes.
\end{enumerate}

Parallel execution is supported via \texttt{ProcessPoolExecutor}, with
\texttt{run\_all} and \texttt{run\_all\_with\_target} orchestrating batch
execution across all cases in a pack.

\subsection{Coverage Analysis}

The \texttt{attacks.coverage} module computes a \texttt{CoverageReport}
from loaded test cases, providing:
\begin{itemize}[nosep]
  \item Per-threat case counts.
  \item Per-strategy case counts.
  \item A threat-strategy matrix showing coverage density.
  \item Gap identification: given expected threats and strategies (from an
        optional \texttt{attacks.yaml} manifest), the module reports missing
        combinations.
\end{itemize}

\subsection{Query Minimization}

The \texttt{attacks.minimize} module implements deterministic delta
debugging (Algorithm~\ref{alg:ddmin}) for failure-preserving query reduction.
Given a failing query and an oracle predicate, the algorithm produces a
minimal subsequence that still triggers the failure.

\begin{algorithm}[H]
\caption{Deterministic Delta Debugging (\texttt{ddmin})}
\label{alg:ddmin}
\begin{algorithmic}[1]
\Require Sequence $c$ of tokens, oracle $\phi: \text{seq} \to \{0,1\}$
\Ensure Minimal $c' \subseteq c$ s.t.\ $\phi(c') = 1$
\State $n \gets 2$
\While{$|c| \geq 2$}
  \State Split $c$ into $n$ disjoint chunks $\Delta_1, \ldots, \Delta_n$
  \State $\mathit{reduced} \gets \mathsf{false}$
  \For{$i = 1$ to $n$}
    \If{$\phi(\Delta_i) = 1$}
      \State $c \gets \Delta_i$; $n \gets 2$; $\mathit{reduced} \gets \mathsf{true}$; \textbf{break}
    \EndIf
  \EndFor
  \If{not $\mathit{reduced}$}
    \For{$i = 1$ to $n$}
      \State $c_i^{\complement} \gets c \setminus \Delta_i$
      \If{$\phi(c_i^{\complement}) = 1$}
        \State $c \gets c_i^{\complement}$; $n \gets n-1$; $\mathit{reduced} \gets \mathsf{true}$; \textbf{break}
      \EndIf
    \EndFor
  \EndIf
  \If{not $\mathit{reduced}$}
    \If{$n < |c|$}
      \State $n \gets \min(2n, |c|)$
    \Else
      \State \textbf{break}
    \EndIf
  \EndIf
\EndWhile
\State \Return $c$
\end{algorithmic}
\end{algorithm}

The splitting strategy is configurable: sentence-based (splitting on
sentence boundaries) or line-based (splitting on newline characters).

% ══════════════════════════════════════════════════════════════════════
\section{Leakage and Integrity Metrics}\label{sec:metrics}
% ══════════════════════════════════════════════════════════════════════

All metrics are implemented as pure functions that take run artifacts and
corpus data as input and produce typed result objects. No metric requires
neural inference, GPU resources, or network access.

\subsection{Canary Detection}

\begin{equation}
\texttt{detect\_canary}(a) = \bigl(\,
  \mathbb{1}[\texttt{CANARY\_[0-9a-f]\{16,\}} \in a],\;
  |\text{matches}|
\,\bigr)
\end{equation}

Returns a \texttt{CanaryResult} with boolean \texttt{present}, integer
\texttt{count}, and list of matched token strings.

\subsection{Verbatim Overlap}

For answer $a$ and source documents $S = \{(d_i, t_i)\}$:

\begin{equation}
\texttt{verbatim\_score} = \max_{(d_i, t_i) \in S} \max\!\left(
  \frac{|\texttt{LCS}(a, t_i)|}{\max(|t_i|, 1)},\;
  \frac{|\texttt{ngram\_matches}(a, t_i, n)|}{|\texttt{ngrams}(t_i, n)|}
\right)
\end{equation}

where $\texttt{LCS}$ returns the longest common substring (minimum length
threshold applied), and n-gram matching uses $n = 5$ by default. The result
(\texttt{VerbatimResult}) includes the score, LCS length, source document
with maximum overlap, and n-gram match count.

\subsection{Membership Confidence}

\textbf{Heuristic approach.}
Given a list of \texttt{RunArtifact} objects, the membership score aggregates
normalized retrieval scores (mean and max across chunks), answer length
(log-normalized), and context specificity (ratio of unique tokens).
When non-member reference data is available, the score is the difference
between member and non-member signal strengths, normalized to $[0, 1]$.

\textbf{Consistency approach.}
For query $q$, generate $k$ paraphrases $\{q_1, \ldots, q_k\}$ using
fixed templates. For each $q_i$, retrieve documents $R_i$ and generate
answer $a_i$. The consistency score combines:

\begin{equation}
\texttt{retrieval\_consistency} = \frac{1}{\binom{k}{2}}
\sum_{i < j} \frac{|R_i \cap R_j|}{|R_i \cup R_j|}
\end{equation}

\begin{equation}
\texttt{answer\_similarity} = \frac{1}{\binom{k}{2}}
\sum_{i < j} \frac{|\text{tokens}(a_i) \cap \text{tokens}(a_j)|}{|\text{tokens}(a_i) \cup \text{tokens}(a_j)|}
\end{equation}

\begin{equation}
\texttt{confidence} = w_r \cdot \texttt{retrieval\_consistency}
                     + w_a \cdot \texttt{answer\_similarity}
\end{equation}

where $w_r = 0.6$ and $w_a = 0.4$ are fixed weights.

\subsection{Semantic Leakage}

For answer $a$ and claims $\mathcal{K}_R$ from retrieved documents:

\begin{equation}
\texttt{semantic\_leakage\_rate} = \frac{|\{c \in \mathcal{K}_R \mid \texttt{match\_claim}(c, a)\}|}{|\mathcal{K}_R|}
\end{equation}

The \texttt{match\_claim} function dispatches to type-specific matchers
(Section~\ref{subsec:semantic}). Text normalization includes whitespace
collapse, punctuation stripping, and case folding.

\subsection{Cross-Document Leakage}

For composed claims $\mathcal{K}_{\times}$ (each spanning $\geq 2$ documents):

\begin{equation}
\texttt{crossdoc\_leaked} = |\{c \in \mathcal{K}_{\times} \mid
  \texttt{match\_claim}(c, a) \;\land\;
  |\texttt{source\_docs}(c)| \geq 2\}|
\end{equation}

\subsection{Integrity Metrics}

Three metrics evaluate corpus poisoning effects:

\begin{equation}
\texttt{poison\_rate@}k = \frac{|\{d \in \text{top-}k \mid d \in \mathcal{P}\}|}{k}
\end{equation}

\begin{equation}
\texttt{poison\_mrr} = \frac{1}{\min_{d \in \mathcal{P}} \text{rank}(d)}
\end{equation}

\begin{equation}
\texttt{true\_doc\_recall@}k = \frac{|\{d \in \text{top-}k \mid d \in \mathcal{T}\}|}{|\mathcal{T}|}
\end{equation}

where $\mathcal{P}$ is the set of poisoned document IDs, $\mathcal{T}$ is
the set of expected true document IDs, and the top-$k$ set is the retrieval result.

\subsection{Verdict Engine}

The \texttt{metrics.verdict} module applies configurable thresholds to
aggregate metric results and produces a final pass/fail verdict.
Default thresholds are:

\begin{itemize}[nosep]
  \item Canary: any extraction $\Rightarrow$ FAIL.
  \item Verbatim: score $> 0.1 \Rightarrow$ FAIL.
  \item Membership: confidence $> 0.7 \Rightarrow$ FAIL.
\end{itemize}

Each threshold violation generates a \texttt{VerdictReason} with the
threat class, human-readable description, observed value, and configured
threshold.

% ══════════════════════════════════════════════════════════════════════
\section{Dataset Poisoning Evaluation}\label{sec:poisoning}
% ══════════════════════════════════════════════════════════════════════

Corpus poisoning represents a supply-chain attack on RAG systems. An
adversary with write access to a subset of the corpus can inject documents
designed to manipulate retrieval ranking, corrupt generated claims, or
implant backdoor triggers.

\subsection{Poisoning Pack Structure}

Each poisoning evaluation is organized as a \emph{pack} containing:
\begin{itemize}[nosep]
  \item A manifest (\texttt{manifest.yaml}) defining queries, expected
        document IDs, expected claims, and poison document identifiers.
  \item A clean corpus and a poisoned corpus variant (produced via the
        \texttt{corpus.patch} mechanism).
  \item Rule definitions for sentinel detection (pattern/replacement pairs
        with deterministic block/strip/allow policies).
\end{itemize}

\subsection{Evidence Collection}

The poisoning runner executes the same queries against both clean and
poisoned corpus variants, collecting \texttt{IntegrityEvidence} objects:

\textbf{Retrieval integrity.} Compares top-$k$ retrieval results between
clean and poisoned runs to detect ranking manipulation. Evidence includes
expected vs.\ actual document ID lists, confidence score, and
per-query metrics (poison rate, MRR, recall).

\textbf{Claim integrity.} Compares generated output against expected true
claims and known poison claims to detect content corruption. Evidence
includes matched true claims, matched poison claims, contradiction count,
and per-query corruption rate.

\textbf{Sentinel integrity.} Applies rule-based patterns to detect
backdoor trigger activation. Evidence includes sentinel type
(suffix/trigger/backdoor), whether the trigger fired, the policy action
taken, expected vs.\ actual behavior, and output markers found.

\subsection{Integrity Summary}

The \texttt{IntegritySection} model aggregates all evidence into a
summary with severity-stratified counts (high/medium/low), per-type
counts (retrieval/claim/sentinel), and worst-case dominance scores.
Evidence is sorted deterministically by severity, then pack ID, then
query ID.

\subsection{Triage}

The integrity summary provides three triage signals:
\begin{itemize}[nosep]
  \item \texttt{worst\_poison\_dominance}: highest \texttt{poison\_rate@k}
        across all retrieval findings.
  \item \texttt{worst\_claim\_corruption}: highest poison claim rate across
        all claim findings.
  \item \texttt{sentinel\_failures}: count of triggered sentinel findings.
\end{itemize}

These signals are consumed by CI gates to determine whether poisoning
exceeds acceptable thresholds.

% ══════════════════════════════════════════════════════════════════════
\section{Calibration}\label{sec:calibration}
% ══════════════════════════════════════════════════════════════════════

Threshold selection for binary metrics (e.g., verbatim score, membership
confidence) requires balancing sensitivity against false positive rate.
RAGLeakLab provides a principled calibration workflow.

\subsection{Labeled Calibration Sets}

Calibration requires a JSONL file of labeled test cases:
\begin{lstlisting}
{"test_id": "case_001", "label": "positive"}
{"test_id": "case_002", "label": "negative"}
\end{lstlisting}

Labels are \texttt{positive} (leak expected, attack succeeds) or
\texttt{negative} (no leak expected, security holds). The loader
(\texttt{calibration.loader}) validates uniqueness, format, and
non-emptiness.

\subsection{Threshold Fitting}

The \texttt{calibration.fit.fit\_threshold} function implements a
sweep-based algorithm:

\begin{enumerate}[nosep]
  \item Separate scores by label into positive and negative sets.
  \item Enumerate all unique score values as candidate thresholds.
  \item For each threshold $t$, compute FPR and TPR:
    \begin{align}
    \text{FPR}(t) &= \frac{|\{s \in S^{-} \mid s \geq t\}|}{|S^{-}|} \\
    \text{TPR}(t) &= \frac{|\{s \in S^{+} \mid s \geq t\}|}{|S^{+}|}
    \end{align}
    where $S^{+}$ and $S^{-}$ are positive and negative score sets
    (assuming \texttt{higher\_is\_worse = True}).
  \item Select the threshold with maximum TPR subject to
    $\text{FPR} \leq \text{target\_fpr}$ (default: 0.01).
  \item Deterministic tie-breaking: among thresholds achieving equal
    TPR at the constraint boundary, prefer the higher (more conservative)
    threshold.
\end{enumerate}

The result (\texttt{CalibrationResult}) reports the chosen threshold,
achieved FPR and TPR, sample counts, and the decision rule string
(e.g., ``score $\geq$ threshold $\to$ FAIL'').

\subsection{Calibration Report}

The \texttt{calibration.report} module generates a full report including:
\begin{itemize}[nosep]
  \item Pack name and metric name.
  \item Target FPR and achieved rates.
  \item ROC-like table of (threshold, FPR, TPR) operating points.
  \item Generation timestamp.
\end{itemize}

Reports are written as JSON to support automated threshold update
workflows (e.g., a CI job that recalibrates thresholds when the corpus
changes and commits updated baseline files).

% ══════════════════════════════════════════════════════════════════════
\section{Attribution}\label{sec:attribution}
% ══════════════════════════════════════════════════════════════════════

When a security test fails, practitioners need to understand \emph{why}
the leak occurred to implement effective remediation. The
\texttt{analysis.attribution} module provides rule-based root cause
analysis.

\subsection{Attribution Categories}

Five categories are implemented, each with a corresponding remediation hint:

\begin{table}[H]
\centering
\caption{Attribution categories and remediation hints.}
\label{tab:attribution}
\small
\begin{tabular}{@{}p{3.2cm}p{4.2cm}@{}}
\toprule
\textbf{Category} & \textbf{Remediation Hint} \\
\midrule
\texttt{retrieval\_included\_secret} &
  Review retriever filtering; consider excluding documents with sensitive markers. \\
\texttt{context\_too\_long} &
  Reduce context window or implement summarization. \\
\texttt{top\_k\_too\_high} &
  Lower top\_k to reduce attack surface. \\
\texttt{chunking\_boundary} &
  Adjust chunk boundaries to avoid splitting sensitive content. \\
\texttt{target\_overexposed\_endpoint} &
  Audit HTTP target for unintended data exposure. \\
\bottomrule
\end{tabular}
\end{table}

\subsection{Attribution Logic}

The \texttt{attribute\_leak} function evaluates conditions based on
run artifact data:
\begin{itemize}[nosep]
  \item If canary detected and retrieved chunks are non-empty, attribute
        to \texttt{retrieval\_included\_secret}.
  \item If context exceeds 10,000 characters, attribute to
        \texttt{context\_too\_long}.
  \item If more than 5 chunks were retrieved, attribute to
        \texttt{top\_k\_too\_high}.
  \item If the target is HTTP-based and a leak was detected, attribute to
        \texttt{target\_overexposed\_endpoint}.
\end{itemize}

Attribution results are embedded in SARIF output as properties on each
finding, enabling security dashboard integration.

% ══════════════════════════════════════════════════════════════════════
\section{CI Integration}\label{sec:ci}
% ══════════════════════════════════════════════════════════════════════

RAGLeakLab is designed as a CI gate: a process that runs in the pipeline,
produces a verdict, and fails the build on security regressions.

\subsection{Export Formats}

\textbf{JUnit XML.} The \texttt{reporting.export.export\_junit} function
produces standard JUnit XML consumed by all major CI platforms (GitHub
Actions, GitLab CI, Jenkins, Azure DevOps). Each test case is mapped to
a \texttt{<testcase>} element; failures include the threat type, query,
and truncated answer. Suppressed findings are marked as
\texttt{<skipped>} with the suppression reason.

\textbf{SARIF.} The \texttt{reporting.export.export\_sarif} function
produces SARIF 2.1.0 for GitHub Code Scanning. Rules are defined for
each threat class with calibrated severity scores:
\begin{itemize}[nosep]
  \item Canary extraction: severity 9.0 (error).
  \item Semantic leakage: severity 8.0 (error).
  \item Cross-document leakage: severity 8.5 (error).
  \item Verbatim leakage: severity 7.0 (warning).
  \item Membership inference: severity 5.0 (warning).
  \item Retrieval hijacking: severity 8.5 (error).
  \item Claim corruption: severity 8.0 (error).
  \item Sentinel takeover: severity 9.0 (error).
\end{itemize}

Suppressed findings receive SARIF \texttt{suppressions} annotations with
\texttt{kind: inSource} and the justification text.

\subsection{HTTP Cassettes}

The cassette mechanism (\texttt{targets.cassette}) enables deterministic
CI runs without network access:

\textbf{Recording.} During a \texttt{record} mode run, each HTTP request
and response is serialized to a JSONL cassette file. Requests are
normalized (sorted headers, sorted query parameters, normalized body) and
hashed (SHA-256) for lookup. Sensitive headers
(\texttt{Authorization}, \texttt{X-Api-Key}, \texttt{Cookie}, etc.)
are redacted before storage.

\textbf{Replay.} During a \texttt{replay} mode run, incoming requests are
normalized and hashed identically. The cassette file is searched for a
matching hash. If found, the stored response is returned without network
access. If not found, a \texttt{CassetteReplayError} is raised.

This design ensures that CI pipelines can execute the full test suite
against recorded external service responses, eliminating network
dependencies and rate-limit concerns.

\subsection{Regression Detection}

The \texttt{regression.diff} module compares two \texttt{ReportSummary}
objects (baseline vs.\ current) and detects regressions:

\begin{itemize}[nosep]
  \item \textbf{Canary regression}: \texttt{false} $\to$ \texttt{true} is
        an immediate failure.
  \item \textbf{Verbatim regression}: rate increase exceeding a
        configurable delta threshold (default: 0.01).
  \item \textbf{Membership regression}: confidence increase exceeding a
        configurable delta threshold (default: 0.05).
  \item \textbf{Semantic regression}: leakage rate increase exceeding the
        verbatim delta threshold.
\end{itemize}

The \texttt{DiffResult} includes per-metric deltas with baseline and
current values, enabling trend analysis.

\subsection{Suppressions}

The suppression system (\texttt{suppressions}) allows teams to acknowledge
known risks without removing them from reports. Suppressions are defined
in \texttt{suppressions.yaml} with mandatory fields:

\begin{itemize}[nosep]
  \item \texttt{id}: UUID identifying the suppression.
  \item \texttt{type}: What is suppressed (\texttt{test\_id},
        \texttt{claim\_id}, \texttt{doc\_id}, or \texttt{metric}).
  \item \texttt{value}: Identifier to match.
  \item \texttt{reason}: Non-empty justification (enforced by validation).
  \item \texttt{expires\_at}: Mandatory ISO-8601 expiry timestamp.
  \item \texttt{owner}: Optional responsible party.
\end{itemize}

Suppressions are applied \emph{after} metric calculation but \emph{before}
the final verdict. Suppressed findings remain visible in all outputs
(JUnit: \texttt{<skipped>}; SARIF: \texttt{suppressions} annotation;
JSON: \texttt{suppression\_summary} section). The verdict changes from
FAIL to KNOWN\_RISK, never to PASS. Expired suppressions are automatically
ignored.

% ══════════════════════════════════════════════════════════════════════
\section{Benchmarks}\label{sec:benchmarks}
% ══════════════════════════════════════════════════════════════════════

RAGLeakLab's design targets two benchmarking dimensions: determinism
verification and execution performance.

\subsection{Determinism Verification}

The \texttt{core.determinism.verify\_determinism} function provides
automated determinism testing. Given a pack name and run count $N$, it:

\begin{enumerate}[nosep]
  \item Creates $N$ fresh \texttt{RAGPipeline} instances and executes
        the pack against each.
  \item Normalizes all outputs by stripping volatile fields
        (\texttt{generated\_at}, \texttt{timings}) and sorting entries by
        \texttt{test\_id}.
  \item Performs deep structural comparison of \texttt{report.json} and
        \texttt{runs.jsonl} across all pairs $(0, i)$ for $i \in [1, N)$.
  \item Reports any differences with path-qualified descriptions (e.g.,
        \texttt{aggregates.verbatim\_leakage\_rate: 0.05 != 0.06}).
\end{enumerate}

Under the reference pipeline with identical corpus and attack inputs, all
runs produce bit-identical normalized outputs. This property is verified in
the test suite.

\subsection{Performance Characteristics}

The framework is designed for lightweight execution:

\begin{itemize}[nosep]
  \item \textbf{No GPU required.} All operations use CPU-only algorithms
        (TF-IDF, regex, string matching, hashing).
  \item \textbf{Parallel execution.} \texttt{run\_all} uses
        \texttt{ProcessPoolExecutor} for multi-core utilization.
  \item \textbf{Caching.} The \texttt{DiskCache} eliminates redundant
        pipeline executions when inputs are unchanged.
  \item \textbf{Cassette replay.} Network-free replay eliminates HTTP
        latency and rate-limit delays.
\end{itemize}

Typical execution times for a 50-case pack against the reference pipeline
on a single CPU core are under 5 seconds, making RAGLeakLab suitable for
pre-commit hooks and PR-blocking CI gates.

% ══════════════════════════════════════════════════════════════════════
\section{Security Hardening}\label{sec:security}
% ══════════════════════════════════════════════════════════════════════

As a security testing tool, RAGLeakLab must not itself become an attack
vector. The following hardening measures are implemented:

\subsection{SSRF Protection}

The \texttt{targets.ssrf} module prevents Server-Side Request Forgery
when using the \texttt{HttpTarget}. The \texttt{validate\_url} function:

\begin{enumerate}[nosep]
  \item Rejects non-HTTP(S) schemes.
  \item Resolves the hostname to its IP addresses.
  \item Blocks connections to private/reserved IP ranges: 10.0.0.0/8,
        172.16.0.0/12, 192.168.0.0/16, 127.0.0.0/8, 169.254.0.0/16,
        ::1/128, fc00::/7, fe80::/10, and the unspecified addresses
        0.0.0.0 and ::.
  \item Optionally enforces a domain allowlist, rejecting any target not
        in the approved set.
\end{enumerate}

\subsection{Secret Redaction}

The \texttt{core.redact} module provides context-aware redaction for
report outputs and logs:

\begin{itemize}[nosep]
  \item Pattern-based redaction: email addresses, phone numbers, canary
        tokens, bearer/basic auth tokens, API keys (Stripe-style, AWS
        AKIA-prefix), and generic credential patterns.
  \item Header-aware redaction: values of sensitive HTTP headers
        (\texttt{Authorization}, \texttt{X-Api-Key}, \texttt{Cookie},
        etc.) are replaced with \texttt{[REDACTED]}.
  \item Recursive operation: \texttt{redact\_dict} traverses nested
        dict/list structures, applying redaction to all string values.
  \item Environment variable scanning: error messages are checked
        against the values of environment variables whose names
        contain sensitive keywords (KEY, SECRET, TOKEN, PASSWORD,
        CREDENTIAL, AUTH).
\end{itemize}

\subsection{Error Message Sanitization}

The \texttt{core.errors} module ensures that error messages displayed to
users or written to CI logs do not leak secrets:

\begin{itemize}[nosep]
  \item The \texttt{sanitize\_message} function applies the same
        redaction patterns used for report output.
  \item Home directory paths are replaced with \texttt{\textasciitilde}
        via \texttt{safe\_path\_str}.
  \item Structured \texttt{SafeError} objects carry sanitized messages
        alongside structured exit codes.
\end{itemize}

\subsection{Filesystem Safety}

The \texttt{core.fs} module provides:

\begin{itemize}[nosep]
  \item \textbf{Path traversal protection}: \texttt{safe\_join} resolves
        paths and verifies the result is within the base directory,
        rejecting absolute paths and \texttt{..} traversal.
  \item \textbf{Atomic writes}: \texttt{atomic\_write} uses a
        temp-file-plus-rename pattern, ensuring files are either fully
        written or untouched. This prevents partial report writes on
        crashes or interrupts.
\end{itemize}

\subsection{Cassette Security}

HTTP cassettes redact sensitive headers before storage and normalize
requests to prevent cache poisoning via header manipulation. Cassette
files use deterministic request hashing (SHA-256 of normalized method,
URL, sorted headers, and body), making it infeasible to craft a request
that maps to a different cassette entry.

% ══════════════════════════════════════════════════════════════════════
\section{Limitations}\label{sec:limitations}
% ══════════════════════════════════════════════════════════════════════

RAGLeakLab makes explicit design trade-offs that result in the following
limitations:

\textbf{No neural-based analysis.} All metrics use string matching,
heuristics, and pattern detection. This ensures determinism and eliminates
GPU dependencies, but means the framework cannot detect paraphrased leakage
that evades substring matching or semantic similarity that requires
embedding-based comparison. The claim matching system uses normalized
substring containment, which may miss reformulations that preserve meaning
while changing surface form.

\textbf{Reference pipeline is not production-representative.} The built-in
TF-IDF retriever and mock generator serve as a controlled baseline, not as
a proxy for production embedding models or language generators. Production
testing requires the \texttt{HttpTarget} adapter or a custom
\texttt{InProcessTarget} wrapping the actual system. The reference pipeline's
value lies in verifying framework behavior, not in predicting production
leakage rates.

\textbf{Fixed paraphrase templates.} The eight built-in paraphrase templates
provide lightweight diversification for consistency-based membership inference.
They do not cover the full space of semantic-preserving reformulations and
may underestimate membership inference risk against systems with more
sophisticated query understanding.

\textbf{No runtime moderation testing.} RAGLeakLab tests the RAG system as
a black box and does not model or test runtime content moderation layers
(e.g., output filters, guardrail classifiers). If a production system
includes such layers, the framework tests the combined system but cannot
attribute filtering to specific moderation components.

\textbf{Claim annotation is manual.} Semantic leakage and cross-document
detection require human-authored claim annotations in JSONL format. The
quality and completeness of these annotations directly bound the metrics'
recall. No automated claim extraction is implemented.

\textbf{Single-process corpus assumption.} The framework assumes
single-process access to the corpus directory during patching and testing.
Concurrent modifications may produce non-deterministic results unless
external locking is applied.

\textbf{Heuristic attribution.} The attribution engine uses rule-based
heuristics with fixed thresholds (e.g., context length $>$ 10,000 characters).
These thresholds are not calibrated per-system and may require tuning for
specific deployment configurations.

% ══════════════════════════════════════════════════════════════════════
\section{Related Work}\label{sec:related}
% ══════════════════════════════════════════════════════════════════════

\textbf{LLM red-teaming frameworks.} Tools such as Microsoft's
Counterfit, NVIDIA's NeMo Guardrails evaluation suite, and various
prompt injection benchmarks focus on adversarial testing of language models.
These operate at the model level and typically require neural inference
for both attack generation and evaluation. RAGLeakLab differs by targeting
the RAG architecture specifically and by eliminating neural dependencies
for deterministic CI operation.

\textbf{Membership inference attacks.} The membership inference literature
originates in the ML privacy domain, with shadow model approaches and
loss-based metrics. RAGLeakLab's heuristic and consistency-based approaches
avoid model training requirements while providing practical signals for
corpus membership detection.

\textbf{Data poisoning.} Backdoor attacks on retrieval systems and
knowledge bases have been studied in the context of search engine
optimization and knowledge graph corruption. RAGLeakLab provides a
structured evaluation framework with three poisoning sub-categories
(retrieval hijacking, claim corruption, sentinel takeover) and
deterministic evidence collection.

\textbf{Canary-based auditing.} The use of planted tokens for
memorization detection was introduced for language model training data
extraction. RAGLeakLab adapts this technique for RAG systems, where the
leakage path is through retrieval and context injection rather than
model memorization.

\textbf{Software testing determinism.} The broader software testing
literature emphasizes flaky test elimination and deterministic execution.
RAGLeakLab applies these principles to security testing, with explicit
volatile field stripping, sorted serialization, and automated
determinism verification.

\textbf{CI security scanning.} Tools like Semgrep, CodeQL, and Snyk
provide static and dependency analysis in CI pipelines. RAGLeakLab
provides dynamic security testing for RAG system behavior, complementing
code-level scanning with runtime leakage detection. The framework's SARIF
and JUnit outputs integrate directly with the same CI platforms.

% ══════════════════════════════════════════════════════════════════════
\section{Conclusion}\label{sec:conclusion}
% ══════════════════════════════════════════════════════════════════════

We have presented RAGLeakLab, a deterministic security testing framework
for Retrieval-Augmented Generation systems. The framework addresses a gap
in the current tooling landscape: the need for reproducible, CI-compatible
security evaluation of RAG architectures without neural inference
dependencies.

The key contributions are:

\begin{enumerate}[nosep]
  \item \textbf{Deterministic design.} A comprehensive architecture where
        every component---from TF-IDF retrieval to delta-debugging query
        minimization---produces bit-identical output across runs.
  \item \textbf{Six threat classes.} Canary extraction, verbatim
        reproduction, membership inference, semantic leakage, cross-document
        correlation, and corpus poisoning, each with well-defined metrics
        and evidence schemas.
  \item \textbf{CI-native integration.} JUnit XML, SARIF 2.1.0, HTTP
        cassettes, regression diffing, and structured suppression
        management enable direct integration into enterprise CI/CD
        pipelines.
  \item \textbf{Principled calibration.} FPR-constrained threshold
        fitting with deterministic tie-breaking enables systematic
        threshold selection for binary metrics.
  \item \textbf{Security hardening.} SSRF protection, secret redaction,
        error sanitization, path traversal prevention, and atomic file
        writes ensure the framework itself does not become an attack
        vector.
\end{enumerate}

RAGLeakLab is open source under the Apache License 2.0 and available at
\url{https://github.com/mishabar410/RAGLeakLab}. The framework requires
only a standard Python installation with no GPU or external service
dependencies for core operation.

Future directions include: integration of lightweight embedding-based
similarity metrics (preserving determinism through pinned model versions),
automated claim extraction from annotated corpora, multi-tenant threat
modeling, and expanded poisoning pack coverage for emerging attack vectors.

\end{document}
